\documentclass[a4paper,11pt]{article}
\usepackage{amsmath,amssymb,amsthm}

\newtheorem{theorem}{Theorem}[section]
\newtheorem{lemma}[theorem]{Lemma}

\newcommand{\abs}[1]{\left\lvert #1 \right\rvert}
\newcommand{\inner}[2]{\left\langle #1, #2 \right\rangle}
\newcommand{\N}{\mathbb{N}}

\title{Notes on Fixed Points and Contractions}
\author{Ren\'e Lefebvre \\ \small Seminar in Analysis, Winter Term}
\date{\today}

\begin{document}
\maketitle

\section{Motivation}

Many problems in analysis reduce to solving an equation of the form $x = T(x)$.
Newton's method, the Picard iteration for differential equations and the
value iteration of dynamic programming are all instances. These notes, typeset
with \LaTeX, collect the basic results; see \S\ref{sec:banach} for the main
theorem and~\cite{banach1922,granas2003} for history and generalisations.

\begin{quote}
  ``The contraction principle is perhaps the most widely applied fixed-point
  theorem in all of analysis.'' --- a remark attributed to many teachers.
\end{quote}

\section{The Banach fixed-point theorem}
\label{sec:banach}

Let $(X, d)$ be a metric space. A map $T \colon X \to X$ is a
\emph{contraction} if there is $q \in [0, 1)$ with
\begin{equation}
  d\bigl(T(x), T(y)\bigr) \le q\, d(x, y) \qquad \text{for all } x, y \in X.
  \label{eq:contraction}
\end{equation}

\begin{lemma}
\label{lem:cauchy}
If $T$ satisfies \eqref{eq:contraction}, then for any $x_0 \in X$ the sequence
$x_{n+1} = T(x_n)$ is Cauchy, and $d(x_m, x_n) \le \frac{q^n}{1-q}\, d(x_1, x_0)$
for $m > n$.
\end{lemma}

\begin{proof}
By induction $d(x_{k+1}, x_k) \le q^k d(x_1, x_0)$. Summing the geometric series
from $k = n$ to $m - 1$ gives the bound, which tends to $0$ as $n \to \infty$.
\end{proof}

\begin{theorem}[Banach, 1922]
\label{thm:banach}
Every contraction of a non-empty complete metric space has exactly one fixed
point, and every Picard sequence converges to it.
\end{theorem}

\begin{proof}
By Lemma~\ref{lem:cauchy} the sequence $(x_n)$ converges to some $x^\ast$, and
continuity of $T$ gives $T(x^\ast) = x^\ast$. If $y^\ast$ were another fixed
point, then $d(x^\ast, y^\ast) \le q\, d(x^\ast, y^\ast)$ forces
$d(x^\ast, y^\ast) = 0$.
\end{proof}

\section{Worked examples}

\begin{enumerate}
  \item The map $T(x) = \cos x$ on $[0, 1]$:
    \begin{enumerate}
      \item $\abs{T'(x)} = \abs{\sin x} \le \sin 1 < 0.85$, so $T$ is a contraction;
      \item the fixed point is the Dottie number, $x^\ast \approx 0.739085$.
    \end{enumerate}
  \item Linear systems $x = Ax + b$ with $\lVert A \rVert < 1$ in some operator norm.
  \item Initial value problems, via the Picard operator
    \[
      (Ty)(t) = y_0 + \int_{t_0}^{t} f\bigl(s, y(s)\bigr)\, ds .
    \]
\end{enumerate}

\begin{center}
  \begin{tabular}{|c|l|r|}
    \hline
    $n$ & $x_n = \cos x_{n-1}$ & error \\
    \hline
    0 & 1.000000 & $2.6 \times 10^{-1}$ \\
    1 & 0.540302 & $2.0 \times 10^{-1}$ \\
    5 & 0.701369 & $3.8 \times 10^{-2}$ \\
    20 & 0.739071 & $1.4 \times 10^{-5}$ \\
    \hline
  \end{tabular}
\end{center}

The convergence is linear with rate $q$, slow compared with Newton's method
but guaranteed from any starting point. Sch\"onberg's and K\"onig's later
refinements are discussed in~\cite{granas2003}; an elementary treatment is
in~\cite{kreyszig1978}.

\begin{thebibliography}{9}
\bibitem{banach1922}
  S.~Banach,
  \emph{Sur les op\'erations dans les ensembles abstraits et leur application aux \'equations int\'egrales},
  Fundamenta Mathematicae \textbf{3} (1922), 133--181.

\bibitem{granas2003}
  A.~Granas and J.~Dugundji,
  \emph{Fixed Point Theory},
  Springer Monographs in Mathematics, Springer, New York, 2003.

\bibitem{kreyszig1978}
  E.~Kreyszig,
  \emph{Introductory Functional Analysis with Applications},
  Wiley, 1978.
\end{thebibliography}

\end{document}
